\href{https://www.llnl.gov/}{\textbf{Lawrence Livermore National Laboratory}},
Livermore, CA, USA

\href{https://computation.llnl.gov/casc/}{\textbf{Center for Applied Scientific Computing}}.
Supervisor:
\href{https://people.llnl.gov/mohror1}{Dr. Kathryn Mohror}

\vspace{-0.5em}

\begin{outerlist}

  \vspace{-0.5em}

  \item[] \textbf{\textit{Graph Science and Irregular Computing Group Leader}}\hfill \textbf{10/2025 -- present}

  \vspace{-0.5em}

  First line manager of a team of 11 computer scientists and applied mathematicians whose work includes the full scope of graph, hypergraph, and network science from theory to advanced software as well as the development of state-of-the-art high performance computing software for managing the irregular communication workloads that arise in graph and data science workloads.
  Responsibilities include advising staff on career planning, connecting them with funding opportunities, and evaluating their performance for the division.
  Additional responsibilities include working with the Computing Directorate's Workforce Division on culture improvement initiatives.
  These include the Thrive Conversations -- a campaign every six months involving interviewing a \textasciitilde{}100-person sample of the \textasciitilde{}1500 Computing employees to identify workforce attitudes towards the Lab culture and any pain points, which we use to deliver anonymized, aggregate guidance to Computing directorate leadership.
  I am an inaugural member of the Thrive Conversation Committee and have been involved since its inception in 2024.
  I am also involved in initiatives intending to address concerns raised in the Thrive Conversations, such as the incipient Stress and Burnout Initiative intended to understand and ameliorate sources of stress among the Computing workforce.

  \vspace{-0.5em}

  \item[] \textbf{\textit{Computing Scientist}}\hfill \textbf{02/2021 -- present}

  \vspace{-0.5em}

  PI and Co-I of multiple research projects investigating scalable graph
  analytics, machine learning, and statistical modeling on High-Performance
  Computing (HPC) systems.
  Project funding sources include
  \href{https://st.llnl.gov/research/internal-investments/ldrd}{LLNL Laboratory Directed Research and Development},
  \href{https://science.osti.gov/ascr}{DOE Office of Science Advanced Scientific Computing Research}, and U.S. Government Agency sponsors.
  Supervised 3 postdocs and dozens of graduate students.
  Selected research contributions include distributed subspace embedding and sketches~\cite{priest2020scaling,krowkee}, distributed K nearest neighbors~\cite{saltatlas,iwabuchi2024neo,iwabuchi2023towards,iwabuchi2026solanet}, distributed power iteration of large structured matrices~\cite{powersqueeze}, distributed clustering of large and high-dimensional data~\cite{clams}, communication infrastructure for asynchronously handling the irregular communication patterns endemic to graph algorithms and related computation~\cite{steil2023embracing,ygm}, novel algorithms and software for scalable Gaussian process (GP) estimation~\cite{muyskens2021muygps,muygpys}, and cosmology, climate, and space domain modeling~\cite{muyskens2022star}.

  \vspace{-0.5em}

  \item[] \textbf{\textit{Data Science Summer Internship Lead}}\hfill \textbf{10/2024 -- present}

  \vspace{-0.5em}

  Manager and facilitator of the
  \href{https://data-science.llnl.gov/dssi}{Data Science Summer Internship},
  a yearly summer program serving \textasciitilde{}40 students and matching LLNL staff mentors each year.
  Responsibilities include managing intern participation in training and networking events such as the
  \href{https://data-science.llnl.gov/dsc}{Data Science Challenge} and the
  \href{https://engineering.llnl.gov/centers/casis/workshops}{CASIS Annual Workshop}, working with mentors to scope and approve summer projects for each student, managing relationships with external partners such as
  \href{https://www.jst.go.jp/EN/}{Japan Science and Technology Agency} and
  \href{https://www.amed.go.jp/en/}{Japan Agency for Medical Research and Development},
  and planning and implementing staff- and professor-taught short courses and seminars that are broadcast to the whole Lab population.

  \vspace{-0.5em}

  \item[] \textbf{\textit{Postdoctoral Researcher}}%
  \hfill \textbf{04/2019 -- 02/2021}

  \vspace{-0.5em}

  Developed novel sketching algorithms to cluster~\cite{priest2020scaling} and
  perform local query approximation~\cite{priest2020degreesketch} massive graphs
  on HPC.
  Solved reinforcement learning~\cite{goumiri2020reinforcement}, image
  classification~\cite{goumiri2020star}, and quantum machine learning~\cite{otten2020quantum} problems using Gaussian Processes (GPs) and neural
  kernels.
  Created an HPC library implementing scalable GP algorithms~\cite{muygpys}.

  \vspace{-0.5em}

  \item[] \textbf{\textit{Computation Student Intern}}%
  \hfill \textbf{05/2018 -- 01/2019}

  \vspace{-0.5em}

  Co-created novel communication library for distributed problems with irregular
  communication patterns, such as sparse linear algebra and graph
  problems~\cite{priest2019you, ygm}.
  Used cardinality sketches to estimate local triangle counts in distributed
  graphs~\cite{priest2018estimating}.

\end{outerlist}
